 %%%%%%%%%%%%%%%%%%%%%%%%%%%%%%%%%%%%%%%%%%%%%%%%%%%%%%%%%%%%%%%%%%%%%%%%%%%%
% AGUJournalTemplate.tex: this template file is for articles formatted with LaTeX
%
% This file includes commands and instructions
% given in the order necessary to produce a final output that will
% satisfy AGU requirements, including customized APA reference formatting.
%
% You may copy this file and give it your
% article name, and enter your text.
%
% guidelines and troubleshooting are here: 

%% To submit your paper:
\documentclass{agujournal2019}
\usepackage{url} %this package should fix any errors with URLs in refs.
\usepackage{lineno}
\usepackage[inline]{trackchanges} %for better track changes. finalnew option will compile document with changes incorporated.
\usepackage{soul}
\usepackage{amsmath}
\linenumbers
%%%%%%%
% As of 2018 we recommend use of the TrackChanges package to mark revisions.
% The trackchanges package adds five new LaTeX commands:
%
%  \note[editor]{The note}
%  \annote[editor]{Text to annotate}{The note}
%  \add[editor]{Text to add}
%  \remove[editor]{Text to remove}
%  \change[editor]{Text to remove}{Text to add}
%
% complete documentation is here: http://trackchanges.sourceforge.net/
%%%%%%%

\draftfalse

%% Enter journal name below.
%% Choose from this list of Journals:
%
% JGR: Atmospheres
% JGR: Biogeosciences
% JGR: Earth Surface
% JGR: Oceans
% JGR: Planets
% JGR: Solid Earth
% JGR: Space Physics
% Global Biogeochemical Cycles
% Geophysical Research Letters
% Paleoceanography and Paleoclimatology
% Radio Science
% Reviews of Geophysics
% Tectonics
% Space Weather
% Water Resources Research
% Geochemistry, Geophysics, Geosystems
% Journal of Advances in Modeling Earth Systems (JAMES)
% Earth's Future
% Earth and Space Science
% Geohealth
%
% ie, \journalname{Water Resources Research}

\journalname{Water Resources Research}


\begin{document}

%%%%%%%%%%%%%%%%%%%%%%%%%%%%%%%%%%%%%%%%%%%%%%%
%  TITLE
%
% (A title should be specific, informative, and brief. Use
% abbreviations only if they are defined in the abstract. Titles that
% start with general keywords then specific terms are optimized in
% searches)
%
%%%%%%%%%%%%%%%%%%%%%%%%%%%%%%%%%%%%%%%%%%%%%%%

% Example: \title{This is a test title}

\title{A comprehensive calibration approach for the Northwest River Forecast Center}

%%%%%%%%%%%%%%%%%%%%%%%%%%%%%%%%%%%%%%%%%%%%%%%
%
%  AUTHORS AND AFFILIATIONS
%
%%%%%%%%%%%%%%%%%%%%%%%%%%%%%%%%%%%%%%%%%%%%%%%

% Authors are individuals who have significantly contributed to the
% research and preparation of the article. Group authors are allowed, if
% each author in the group is separately identified in an appendix.)

% List authors by first name or initial followed by last name and
% separated by commas. Use \affil{} to number affiliations, and
% \thanks{} for author notes.
% Additional author notes should be indicated with \thanks{} (for
% example, for current addresses).

% Example: \authors{A. B. Author\affil{1}\thanks{Current address, Antartica}, B. C. Author\affil{2,3}, and D. E.
% Author\affil{3,4}\thanks{Also funded by Monsanto.}}

\authors{G. Walters\affil{1} and C. Bracken\affil{2}}


\affiliation{1}{Northwest River Forecast Center, National Oceanic and Atmospheric Administration, Portland, Oregon, USA}
\affiliation{2}{Pacific Northwest National Laboratory, Richland, Washington, USA}
% \affiliation{3}{Third Affiliation}
% \affiliation{4}{Fourth Affiliation}

%(repeat as many times as is necessary)


% Corresponding author mailing address and e-mail address:

% (include name and email addresses of the corresponding author. More
% than one corresponding author is allowed in this LaTeX file and for
% publication; but only one corresponding author is allowed in our
% editorial system.)

% Example: \correspondingauthor{First and Last Name}{email@address.edu}

\correspondingauthor{Geoffrey Walters}{geoffrey.walters@noaa.gov}



%%%%%%%%%%%%%%%%%%%%%%%%%%%%%%%%%%%%%%%%%%%%%%%
% KEY POINTS
%%%%%%%%%%%%%%%%%%%%%%%%%%%%%%%%%%%%%%%%%%%%%%%
%  List up to three key points (at least one is required)
%  Key Points summarize the main points and conclusions of the article
%  Each must be 140 characters or fewer with no special characters or punctuation and must be complete sentences

% Example:
% \begin{keypoints}
% \item	List up to three key points (at least one is required)
% \item	Key Points summarize the main points and conclusions of the article
% \item	Each must be 140 characters or fewer with no special characters or punctuation and must be complete sentences
% \end{keypoints}

\begin{keypoints}
\item enter point 1 here
\item enter point 2 here
\item enter point 3 here
\end{keypoints}

%%%%%%%%%%%%%%%%%%%%%%%%%%%%%%%%%%%%%%%%%%%%%%%
%
%  ABSTRACT and PLAIN LANGUAGE SUMMARY
%
% A good Abstract will begin with a short description of the problem
% being addressed, briefly describe the new data or analyses, then
% briefly states the main conclusion(s) and how they are supported and
% uncertainties.

% The Plain Language Summary should be written for a broad audience,
% including journalists and the science-interested public, that will not have 
% a background in your field.
%
% A Plain Language Summary is required in GRL, JGR: Planets, JGR: Biogeosciences,
% JGR: Oceans, G-Cubed, Reviews of Geophysics, and JAMES.
% see http://sharingscience.agu.org/creating-plain-language-summary/)
%
%%%%%%%%%%%%%%%%%%%%%%%%%%%%%%%%%%%%%%%%%%%%%%%

%% \begin{abstract} starts the second page

\begin{abstract}
[ enter your Abstract here ]
\end{abstract}

\section*{Plain Language Summary}
Enter your Plain Language Summary here or delete this section.
Here are instructions on writing a Plain Language Summary: 
\url{https://www.agu.org/Share-and-Advocate/Share/Community/Plain-language-summary}


%%%%%%%%%%%%%%%%%%%%%%%%%%%%%%%%%%%%%%%%%%%%%%%
%
%  BODY TEXT
%
%%%%%%%%%%%%%%%%%%%%%%%%%%%%%%%%%%%%%%%%%%%%%%%

%%% Suggested section heads:
% \section{Introduction}
%
% The main text should start with an introduction. Except for short
% manuscripts (such as comments and replies), the text should be divided
% into sections, each with its own heading.

% Headings should be sentence fragments and do not begin with a
% lowercase letter or number. Examples of good headings are:

% \section{Materials and Methods}
% Here is text on Materials and Methods.
%
% \subsection{A descriptive heading about methods}
% More about Methods.
%
% \section{Data} (Or section title might be a descriptive heading about data)
%
% \section{Results} (Or section title might be a descriptive heading about the
% results)
%
% \section{Conclusions}

\section{Introduction}
%Text here ===>>>
\begin{itemize}
    \item History of calibration at NWRFC
    \item Models used
    \item Other calibration studies 
\end{itemize}

\begin{figure}
    \centering
    \includegraphics[width=\textwidth]{figures/basin-map.png}
    \caption{Map showing the NWRFC domain (thick black line), CAMELS basins (filled in grey) and three representative basins used as calibration examples (filled in color). Other basins modeled my the NWRFC are shown with thin grey lines.}
    \label{fig:map}
\end{figure}
\section {Existing Modeling System}

The sections below describe the general lumped rainfall runoff suite of models used by a river forecast center (RFC) to translate precipitation and snow melt to a hydrograph and route downstream. 

A single instance of the following suite of model can be used to represent the entire watershed, or alternatively a collections of model suites can be used to model discrete portion of the watershed, often referred by a RFC as zones. Each zone has it's own unique parameter set for the suite of hydrologic models. 

\subsection{Snow}

The SNOW-17 is a temperature indexed snow accumulation and ablation model. The general physical process of snow cover are captured by the model but in a simplified form. The model uses air temperature and precipitation rates to inform the amount to accumulation of snow melt. 
 Moving through the winter, spring, and into the summer the model becomes more sensitive to temperature and its impact on the meltrate. 
 There is also an accounting component within the model that tracks the snowpacks heat deficit, liquid to water ratio, and snow cover of the area being modeled. The heat deficit needs to rise to 32 degF and the melted water stored in the snowpack has to rise above capacity to begin the onset of melt. The model's state of the snow cover controls the amount of melt, where more more water is made available as melt when there is more simulated snow cover. 

\subsubsection{Glacial Zones}

For glacial zones the the initial snow depth is set at an infinite height and the snow is always assumed to cover the entire area \cite{anderson_calb_manual2002}.

 \begin{figure}[ht]
    \centering
    \includegraphics[width=\textwidth]{figures/Snow17_Diagram.png}
    \caption{Snow17 Diagram}
    \label{fig:cyclical}
\end{figure}
 
\subsection{Soil}

The Sacramento Soil Moisture Accounting Model (SAC-SMA) is a conceptual soil model. The model contain several parameters, each of which represents a discrete portion of the soil runoff process such as impervious runoff, surface runoff, interflow, and baseflow. It also has an soil moisture accounting function to track water moving in, out, and through the soil moisture column via evapotranspiration, soil-moisture storage, and drainage processes. Ultimately, the SAC-SMA model regulates the timing of water travels through the soil column. 
The model parameters are not calibrated by sampling the soil profile via field studies, but rather through inference using observations of rainfall and streamflow runoff. 

\begin{figure}[ht]
    \centering
    \includegraphics[width=\textwidth]{figures/sac-sma diagram.png}
    \caption{SAC-SMA Diagram}
    \label{fig:cyclical}
\end{figure}

Good overview of SACSMA model to pull from:  \cite{sacsma_manual2002}
Original Report:
A Generalized Streamflow Simulation System: Conceptual Modeling for Digital Computers, Robert J.C. Burnash, R. Larry Ferral, Richard A. McGuire, 1973, NOAA/State of California, Report

\subsubsection{Glacial Zones}

For glacial zones only certain portions of the SAC-SMA are used, primarily those that control the attenuation of water through the system. The following Sacramento model components are turned off; tension water storage, impervious flow, evapotranspiration, losses from riparian vegetation. \cite{anderson_calb_manual2002}.

\subsection{Unit Hydrograph}

The unit hydrograph (UH)is a scaled time distribution of a unit depth of excess water over a given area from a rainfall-runoff model. Since SAC-SMA handle the timing of runoff through the soil column, the NWRFS unit hydrograph does not consider the excess water until it enters the channel system. This is different from traditional unit hydrograph approaches which consider interflow travel paths through the soil column as part of the unit hydrograph shape,   

okay overview of UH:
\cite{unithg_manual2002}
Original UH document
\cite{linsley1982hydrology}

\subsection{Routing}

Typically, RFC route a simulation downstream to the next location being modeled using a hydrologic routing approach. If needed, any flow into the stream channel between the two modeling locations is captured by a collection of additional Snow17, SAC-SMA, and UH suites to represent that local area. There are multiple hydrologic routing models used by RFCs, but for the most recent NWRFC calibration work, the Lag and K routing method (lagk) was exclusively used. It is flexible method of routing since both the Lag and K elements can be either constant or variable by streamflow (Linsley, Kohler and Paulhus, 1975, Section 9.9). The lag parameter generates a shift in the hydrograph forward to represent travel time and the K parameter represents the attenuation of the streamflow as it travels downstream. 

\subsection{Additional modules}

There are two additional hydrologic modules utilized at the NWRFC; Channel Loss and Consumptive Use. The Channel Loss (CHANLOSS) operation accounts for losses or gains of water that occur along a channel reach as a result of interactions with groundwater, anthropogenic impacts, or evaporation from the stream surface. The Consumptive Use (CONS\_USE) operation accounts for the impacts of surface water irrigation on streamflow based on crop evapotranspiration demand and irrigation efficiency.

okay overview of Lagk:
\cite{lagk_manual2002}
Original LagK document
\cite{linsley1982hydrology}

\section{Data}

\subsection{Forcings}

The National Weather Service's Analysis of Record for Calibration (AORC), is a multi-decade, high-resolution dataset containing all weather information necessary for forcing land-surface, snow, and hydrologic models has been developed. It includes the period 1979-Present, at a time interval of one hour. Data are stored on a 30” (0.008333°) latitude/longitude grid mesh. The primary motivation for developing this dataset was the need for a climatology-constrained near-surface weather record suitable for calibrating hydrologic models \cite{Fall_2023}. 

The dataset variables used for this study are 2-m above ground air temperature, specific humidity, terrain-level pressure, and, and 1-h precipitation accumulation. Variables other than precipitation are instantaneous at the start of the hour; precipitation accumulation ends at the given hour. 

For each time step, zone average values were calculated for all AORC varibles used as input into the model. 

\subsection{Streamflow data}

Hydrologic models developed for operational forecasting by a RFC are typically set at a locations which are also occupied  by an established measurement station that is collecting streamflow observations. Prior to beginning a calibration, all daily average streamflow observations are collected at calibration site for same period there was AORC data availability. When daily streamflow was unavailable, for example the station wasn't established until sometime after the start of AORC availability, the data would be imputed. 

Imputation of missing streamflow data was conducted using the missRanger R package \cite{Mayer2024}, an implementation of the MissForest algorithm \cite{Stekhoven_2011}. To assess the quality of the imputed data we conducted a cross-validation analysis using available stream flow records by randomly dropping 50\% of the data at a time and imputing the rest. The imputed values were compared to observations and performance was assessed using correlation ($R^2$), NSE, and percent bias. This analysis was conducted 5 times and the results were averaged to asses overall performance of the imputation (Table in supplemental material). Additionally, imputation was conducted using a linear regression model to provide a performance baseline. Performance is generally very favorable with imputed values showing high $R^2$ and NSE and low bias. On average the random forest imputation agorithm performed x\% better than the linear regression model. 


TODO: format results table, maybe some summary stats 

\section{Model enhancements}

\subsection{Dynamic ET}

Evapotranspiration is used as input to SACSMA model and interacts with the upper and lower tension water bucket. The  tension water bucket is a conceptualized relationship that represents that moisture which is attracted to moisture-deficient soil particles so strongly that it can be removed only by evapotranspiration. How much water is removed from the tension water buckets depends on the evaporation demand and the current contents of the upper and lower tension water buckets.

The Hargreave Samani equation was used to calculate daily potential evapotranspiration \cite{Hargreaves_1982} . 
The equation is a temperature indexed evaporation approach, but also incorporates solar radiation.

\begin{equation}
    PET=C_1(T_a+17.8) R_e/\lambda(T_x-T_n )^{C_2 }
\end{equation}

Tx and Tn are max and min temperatures, respectively, and C1 and C2 are fitting coefficients. A unique calibrated C1 coefficient was assigned to the 15th of each calendar month. For days in between the 15th, the coefficient was linearly interpolated using adjacent months. As is common practice, C2 was set to 0.5. 

Although a temperature indexed approach is simple in comparison to a full energy balance method, it is an improvement over legacy NWRFC methods. Previous operational models used a static daily rate regardless of basin conditions. 

The potential evapotranspiration was converted to evapotranspiration demand using available monthly Normalized Difference Vegetation Index (NDVI) to calculated a potential evapotranspiration adjustment factor (PETadj). An average monthly PETadj factor was calculated for the 15th of each calendar month. For days in between the 15th, the coefficient was linearly interpolated using the adjacent months. 


\subsection{Zone Delineation}

As part of the hydrologic model calibration workflow by a RFC the process involves exploring if it is advantageous to split the basin in to multiple discrete modeling zones. In the Western United states, when snow is prevalent in a basin, two zones will be often be used; one zone representing the higher elevations where snowpack persists through the winter and into the spring and another zone where the snowpack is transient. During past calibrations, an elevations would be selected as the line of demarcation and the basin would be split into a upper and lower zone. For certain scenarios multiple elevations would be used to split it into more than two zones. Also glacial zones are typically modeled as their own zone.

As part of this new calibration approach, an unsupervised grouping technique was used to segregate the basin into unique zone. The process relied on collecting gridded basin characteristics (table II). Portions to be modeled as a glacier were removed from the analysis. All grid were resampled to the same resolution. For a given basin, the hydrologic characteristics associated with a given grid cells were grouped. Using the hydrologic characteristics for each grid cell, K mean clustering was used to divide the basins into a specified amount of zones. Because of the role that orthographic effects play in many of the hydrologic characteristics utilized, the kmean clustering method often followed elevation contours but with additional nuances due such as changes in forest cover, soil type, or basin aspect.

\begin{figure}[ht]
    \centering
    \includegraphics[width=\textwidth]{figures/example_delineation.png}
    \caption{Example basin delineation using kmeans clustering}
    \label{fig:cyclical}
\end{figure}

\begin{table}[ht]
    \centering
    \begin{tabular}{lcc}
        \hline
        Dataset                                 & Data Source  & Reference \\
        \hline
        Average PTPS from Nov-Mar               & AORC         & \\ 
        Average Annual Cumulative Precipitation & AORC         & \\ 
        Topographical Elevation                 & USGS         & \\ 
        Effective Forest Cover                  & NLCD, Canada & \citeA{Beaudoin_2014}\\ 
        Saturated Hydraulic Soil Conductivity   &              & \citeA{Zhang_2018}\\ 
        30day Avg Max Annual SWE depth          & U Arizona    & \citeA{Zeng_2018}\\ 
                                                &              & \citeA{Broxton2019}\\ 
        30day Avg Max Annual SWE depth, Canada  & ERA5-Land    & \citeA{Mu_oz_Sabater_2021}\\
                                                &              & \citeA{CCCS2019}\\
        \hline
    \end{tabular}
    \caption{Caption}
    \label{tab:my_label}
\end{table}

Average PTPS from Nov-Mar:  AORC \cite{Fall_2023}

Average Annual Cumulative Precipitation: AORC \cite{Fall_2023}

Topographical Elevation:  USGS \cite{usgs2023}

Effective Forest Cover:  NLCD, Canada:  \cite{Beaudoin_2014}

Saturated Hydraulic Soil Conductivity:  \cite{Zhang_2018}

30day Avg Max Annual SWE depth:  U Arizona dataset (\cite{Zeng_2018,Broxton2019}) and ERA 5 land in Canada \cite{Mu_oz_Sabater_2021,CCCS2019}.

\subsection{Precipitation Typing}

Historically, to derive the percentage of precipitation falling as snow, NWRFC has relied on a temperature time series assigned to a specific basin elevation, lapse rate, and an area elevation curve. 

For most recent recalibrations, specific humidity and terrain-level pressure from the aorc gridded dataset were used to derive wetbulb temperatures. The Western Region National Weather Service Guidance uses a standardized approach of using wet-bulb 0.5 deg to set the snow level \cite{Cleave2019}. Leveraging that guidance, for a given time step,  a typing grid was created of binary values (0|1) depending the value of wetbulb grid at that same location. A zone averaged time series of typing was calculated for each time step to be used as input to the model. 

\section{Automatic calibration}

This section describes the procedure used for automatic calibration of a single basin which involves selecting a calibration algorithm, calibrated parameters, upper and lower parameter limits, and some special considerations for basins which include snow and routed flow. 

%https://ascelibrary.org/doi/10.1061/%28ASCE%29HE.1943-5584.0000938
%https://agupubs.onlinelibrary.wiley.com/doi/10.1029/2005WR004723
%https://www.sharcnet.ca/my/publications/show/1001

For automatic calibration we used the dynamically dimension search (DDS) algorithm \cite{Tolson2007}. DDS has been recognized as robust tool to perform hydrologic model optimizations. The DDS algorithm is not designed to find global optimum parameter values but instead find reasonable parameter values within a given computational budget (i.e. iteration limit). Initial testing with the original DDS algorithm found that when performing a 10,000 iteration run on a single core, the run time and performance was comparable to other statistical optimizers in a comparable total amount of time, such as particle swarm optimization (PSO) which uses multiple cores. This opened up the possibility that each CPU core could run an independent auto-calibration with the DDS optimizer in parallel, with the best result from all the runs used as the final solution. Running the DDS optimizer independently on multiple cores simultaneously is referred to as an embarrassing parallel approach (EP-DDS). Another parallel DDS (P-DDS) implementation approach was explored by \cite{Tolson2014} that instead of the EP-DDS approach, there is a periodic check in and exchange of information between the runs. Tolson found that there is no clear winner between EP-DDS and P-DDS, but each has a distinct advantage depending on how far the algorithm is into the optimization run. He stated that ``The next step in parallel DDS algorithm development involves merging the two parallel implementations to develop a parallel DDS implementation that combines the best aspects of EP-DDS and P-DDS.'' 

Based on the strategies recommended by the author to utilize each strength of EP-DDS and P-DDS, we developed a technique where the check in between independent parallel optimization process progressed as iterations increased. Initially the parallel processes would only exchange information every 1,000 iterations, to closely resemble a DDS-EP run. As the DDS run progresses, the frequency of check-ins increases from 1,000 to 500 to 100 to finally every 10 iterations. As the check in frequency increases, the DDS run more closely resembles the P-DDS process. The technique was coined an evolving DDS (EDDS) parallelization approach. The approach resulted in improved run times compared to EP-DDS and P-DDS while still maintaining comparative calibration skill to other optimizers \cite{Arsenault2014,Asgari2023}.

https://onlinelibrary.wiley.com/doi/full/10.1111/jawr.12143
https://onlinelibrary.wiley.com/doi/full/10.1111/1752-1688.12402

\subsection{Objective function}
Calibration studies often use a single objective function such as root mean squared error (RMSE) or Nashe-Suttecliffe efficiency (NSE) which is minimized to produce an optimal parameter set (reference). Both metrics rely on squared deviation between observed and simulated values and so put more weight on high flows which can misrepresent of low flows, which can be just as important to water managers in the Pacific Northwest for hydropower, fish management, navigation and recreation. For these reasons we selected an combined objective function which balances low and high flows
$$\text{min} \sum_i \text{NSE}(s_i-o_i)+\text{NSE}(\log(s_i)-\log(s_i))$$
where the $s_i$ is the simulated flow and $o_i$ is the observed flow at timestep i. Note that the model was run at a 6 hour timestep but the objective function was computed for daily values. 

\subsection{Basin clustering for limits}

Parameter limits for optimization were determined based on clustering all available basins into groups which share similar hydrologic properties. The properties used for clustering were (1) Average percent precip as snow, (2) Water year total precipitation, (3) Average annual SWE, (4) basin average elevation, (5) effective forest cover percentage, and (6) Average soil conductivity. 

\subsubsection{Calibration of UH}

The USGS national hydrologic dataset was used to calculated the maximum travel time for excess water to reach a basin's outlet. A velocity field was calculated to determine travel time using a method proposed by \citeA{Maidment1996}. A velocity limit of 0.02 to 2 m/sec was assumed. 

A gamma synthetic unit hydrograph was optimized within the auto calibration routine \cite{Croley_1980}. The UH gamma shape and scale parameters were constrained during optimization so that the unit hydrograph shape length was within 25\% of the calculated maximum travel time. {Need more detail on this.}

For prior NWRFC calibrations duplicate unit hydrographs were used for each zone. With the inclusion of auto calibration, each zone had its own optimized unit hydrograph using the calculated its own calculated maximum travel time to the basin outlet. In practice this improves calibrations by allowing the model to be more flexible. In some cases this exacerbates the equifinality of the solution space and causes something akin to label switching in hidden markov models (reference). We observed in some cases of optimization that one zone serves as the baseflow and the other zone serves as the flashy component, and in a second run of the optimizer, the role of each zone switches. This issue was generally avoided by careful use of parameter limits for each zone. 

\subsubsection{Areal Depletion Curve}

In order to apply SNOW-17 to an area, the model tracks areal extent of the snow cover. That extent is derived by developing a calibrated relationship between the depth of the snow water equivalent and the amount of snow cover in the form of a 10 ordinate lookup table. 

If the values of area depleation table were to be plotted in two dimensional space, it is expected that the resultant line be physically reasonable and be smooth in shape. To ensure these two requirements are met, the following equation was used

$$ax^b+(1-a)x^c$$ 

The parameters a,b, and c were optimized with limits then ensured a physically reasonable result.

\subsubsection{LagK Table}

Similar to the areal depletion curve, two lookup tables were used to characterize the lag and attenuation within the LagK model. For both the lag and attenuation the values were associated with flow of the upstream tributary. Also like the area depletion curve, the values withing the table should be physically realistic and of a smooth shape when plotted in two dimensional space. 

For lag and k values, the equation used $$\text{lag or k table entry}=a(Q-d)^{2}+bQ+c$$ . a, b, c, and d were were optimized. 

For streamflow (Q) associated with the table, the observed historical streamflow observations were used to derived upper and lower limits. 
 Lag limits were set using the river miles between the calibration location and the upstream tributary and assuming river speed could be anywhere from one to seven miles an hour. The lag table was limited so that travel time has to decrease with increasing flow.

 \subsection{Optimized parameters}

Generally, parameters were not optimized when a measurable physical connection within the basin could be made. Examples of this would be area, mean elevation, effective forest cover, and latitude. 

Other parameters were not optimized when they represented an assumed physical process limits. These were exclusively Snow-17 parameters and they included:  maximum negative melt factor, anecedent snow temperature, base temperature for non-rain malefactor, maximum percent of liquid water held in a snowpack, daily melt at the snow-soil interface. 

However, the majority of parameters associated with the SNOW-17 and SAC-SMA cannot be directly linked to measurable basin attributes, and were optimized during the calibration process. 

\begin{table}[ht]
    \centering
    \begin{tabular}{cccc}
    \hline
        SAC-SMA  & SNOW17                & Unit Hydrograph & LagK\\
    \hline
        adimp    & areal depletion curve & shape           & K   \\  
        lzfpm    & mfmax                 & scale           & lag \\  
        lzfsm    & mfmin                 &                 &     \\  
        lzpk     & scf                   &                 &     \\  
        lzsk     & si                    &                 &     \\  
        lztwm    & uadj                  &                 &     \\  
        pctim    &                       &                 &     \\  
        pfree    &                       &                 &     \\  
        rexp     &                       &                 &     \\  
        riva     &                       &                 &     \\  
        uzfwm    &                       &                 &     \\  
        uzk      &                       &                 &     \\  
        uztwm    &                       &                 &     \\  
        zperc    &                       &                 &     \\
    \hline
    \end{tabular}
    \caption{Optimized model pentameters}
    \label{tab:my_label}
\end{table}

Below are some special optimization cases for certain parameters or required calibrated tables:

\subsection{Software wrappers}

Should be able to point to repos for both of these

Work on documentation

\section{Calibration Results - case studies}

In this section we provide a detailed look at the calibration results of three representative basins in the NWRFC domain. The first, Nehalem R near Foss, Oregon (USGS \#14301000, NWS ID FSSO3) is a rain dominated basin with a winter peak with limited snowmelt. The second, the middle fork of the Flathead River near West Glacier, Montana (USGS \#12358500, NWS ID WGCM8) is a snow dominated basin with a summer peak and limited winter rainfall. The third, the Sauk River Near Sauk, WA (USGS \#12189500, NWD ID, SAKW1), which is a basin with routed upstream inflow and both summer snowmelt and winter rain. 

Table 1 shows the NSE values for each cross validation period (denoted cv1 through cv4) which drops 1/4 of the data, approximately 10 years, to compute the calibration objective function for the remaining data. A run was also conducted using the entire period of record (POR) for calibration, this represents the calibration that is used operationally. The table also has rows representing the calibration results with forcing adjustments and without. Notably the NSE values are relatively stable across the CV folds and the POR run, which indicates that the auto-calibration is finding similarly optimal solutions. It is especially encouraging that the POR runs results are generally similar in quality to the CV runs which indicates no major data issues. At SAKW1 and WGCM8, the forcing adjustment does little to improve the overall calibrations but at FSSO3, the forcing adjustments dramatically improve the calibration results, indicating a potential issue with the forcings in this basin. These results indicate that forcing adjustment is not necessary in all basins but that it does not degrade the performance when used. 

% latex table generated in R 4.3.3 by xtable 1.8-4 package
% Fri Apr 19 12:59:37 2024
\begin{table}[ht]
\centering
\begin{tabular}{rllrrrrr}
  \hline
  basin & Forcing Adjustment & CV1 & CV2 & CV3 & CV4 & POR \\ 
  \hline
   FSSO3 & Yes & 0.935 & 0.947 & 0.942 & 0.943 & 0.945 \\ 
   FSSO3 & No  & 0.710 & 0.713 & 0.716 & 0.680 & 0.704 \\ 
   SAKW1 & Yes & 0.853 & 0.845 & 0.853 & 0.841 & 0.856 \\ 
   SAKW1 & No  & 0.840 & 0.841 & 0.838 & 0.807 & 0.837 \\ 
   WGCM8 & Yes & 0.956 & 0.964 & 0.963 & 0.958 & 0.956 \\ 
   WGCM8 & No  & 0.956 & 0.960 & 0.953 & 0.953 & 0.957 \\ 
   \hline
\end{tabular}
\caption{NSE values for the three test basins for each cross validation (CV) period and the entire period of record (POR). Also provided are the results with forcing adjustment and without. }
\end{table}



\begin{figure}[!ht]
    \centering
    \includegraphics[width=\textwidth]{figures/timeseries.pdf}
    \includegraphics[width=\textwidth]{figures/cyclical.pdf}
    \caption{Example calibrated simulations from 2019 to 2022}
    \label{fig:timeseries}
\end{figure}

\begin{figure}[ht]
    \centering
    \includegraphics[width=\textwidth]{figures/peak_flow.pdf}
    \includegraphics[width=\textwidth]{figures/low_flow.pdf}
    \caption{Caption}
    \label{fig:peak-flow}
\end{figure}

\section{Calibration Results - Camels Basins}

To facilitate intermodel comparison we present the results of a 

\begin{figure}[ht]
    \centering
    \includegraphics[width=\textwidth]{figures/budyko.pdf}
    \caption{Caption}
    \label{fig:budyko}
\end{figure}

\begin{itemize}
    \item map of all completed runs summary graphic
    \item heat map of KGE score in budyko space 
    \item introduction 
    \item finish calibration analysis 
\end{itemize}

\section{Discussion}
Calibration Guidance

Gold standard for lumped model calibration in the PNW 

\section{Conclusion}
%%

%  Numbered lines in equations:
%  To add line numbers to lines in equations,
%  \begin{linenomath*}
%  \begin{equation}
%  \end{equation}
%  \end{linenomath*}



%% Enter Figures and Tables near as possible to where they are first mentioned:
%
% DO NOT USE \psfrag or \subfigure commands.
%
% Figure captions go below the figure.
% Acronyms used in figure captions will be spelled out in the final, published version.

% Table titles go above tables;  other caption information
%  should be placed in last line of the table, using
% \multicolumn2l{$^a$ This is a table note.}
% NOTE that there is no difference between table caption and table heading in the final, published version
%
%----------------
% EXAMPLE FIGURES
%
% \begin{figure}
% \includegraphics{example.png}
% \caption{caption}
% \end{figure}
%
% Giving latex a width will help it to scale the figure properly. A simple trick is to use \textwidth. Try this if large figures run off the side of the page.
% \begin{figure}
% \noindent\includegraphics[width=\textwidth]{anothersample.png}
%\caption{caption}
%\label{pngfiguresample}
%\end{figure}
%
%
% If you get an error about an unknown bounding box, try specifying the width and height of the figure with the natwidth and natheight options. This is common when trying to add a PDF figure without pdflatex.
% \begin{figure}
% \noindent\includegraphics[natwidth=800px,natheight=600px]{samplefigure.pdf}
%\caption{caption}
%\label{pdffiguresample}
%\end{figure}
%
%
% PDFLatex does not seem to be able to process EPS figures. You may want to try the epstopdf package.
%

%
% ---------------
% EXAMPLE TABLE
%
% \begin{table}
% \caption{Time of the Transition Between Phase 1 and Phase 2$^{a}$}
% \centering
% \begin{tabular}{l c}
% \hline
%  Run  & Time (min)  \\
% \hline
%   $l1$  & 260   \\
%   $l2$  & 300   \\
%   $l3$  & 340   \\
%   $h1$  & 270   \\
%   $h2$  & 250   \\
%   $h3$  & 380   \\
%   $r1$  & 370   \\
%   $r2$  & 390   \\
% \hline
% \multicolumn{2}{l}{$^{a}$Footnote text here.}
% \end{tabular}
% \end{table}

%%%%%%%%%%%%%%%%%%%%%%%%%%%%%%%%%%%%%%%%%%%%%%%
% SIDEWAYS FIGURES and TABLES
% AGU prefers the use of {sidewaystable} over {landscapetable} as it causes fewer problems.
%
% \begin{sidewaysfigure}
% \includegraphics[width=20pc]{figsamp}
% \caption{caption here}
% \label{newfig}
% \end{sidewaysfigure}
%
%  \begin{sidewaystable}
%  \caption{Caption here}
% \label{tab:signif_gap_clos}
%  \begin{tabular}{ccc}
% one&two&three\\
% four&five&six
%  \end{tabular}
%  \end{sidewaystable}

%% If using numbered lines, please surround equations with \begin{linenomath*}...\end{linenomath*}
%\begin{linenomath*}
%\begin{equation}
%y|{f} \sim g(m, \sigma),
%\end{equation}
%\end{linenomath*}

%%% End of body of article

%%%%%%%%%%%%%%%%%%%%%%%%%%%%%%%%%%%%%%%%%%%%%%%
%% Optional Appendices go here
%
% The \appendix command resets counters and redefines section heads
%
% After typing \appendix
%
%\section{Here Is Appendix Title}
% will show
% A: Here Is Appendix Title
%
%\appendix
%\section{Here is a sample appendix}

%%%%%%%%%%%%%%%%%%%%%%%%%%%%%%%%%%%%%%%%%%%%%%%
% Optional Glossary, Notation or Acronym section goes here:
%
% Glossary is only allowed in Reviews of Geophysics
%  \begin{glossary}
%  \term{Term}
%   Term Definition here
%  \term{Term}
%   Term Definition here
%  \term{Term}
%   Term Definition here
%  \end{glossary}


%%%%%%%%%%%%%%%%%%%%%%%%%%%%%%%%%%%%%%%%%%%%%%%
% Acronyms
%% NOTE that acronyms in the final published version will be spelled out when used in figure captions.
%   \begin{acronyms}
%   \acro{Acronym}
%   Definition here
%   \acro{EMOS}
%   Ensemble model output statistics
%   \acro{ECMWF}
%   Centre for Medium-Range Weather Forecasts
%   \end{acronyms}


%%%%%%%%%%%%%%%%%%%%%%%%%%%%%%%%%%%%%%%%%%%%%%%
% Notation
%   \begin{notation}
%   \notation{$a+b$} Notation Definition here
%   \notation{$e=mc^2$}
%   Equation in German-born physicist Albert Einstein's theory of special
%  relativity that showed that the increased relativistic mass ($m$) of a
%  body comes from the energy of motion of the body—that is, its kinetic
%  energy ($E$)—divided by the speed of light squared ($c^2$).
%   \end{notation}




%%%%%%%%%%%%%%%%%%%%%%%%%%%%%%%%%%%%%%%%%%%%%%%
%
% DATA SECTION and ACKNOWLEDGMENTS
%
%%%%%%%%%%%%%%%%%%%%%%%%%%%%%%%%%%%%%%%%%%%%%%%

\section*{Open Research Section}
This section MUST contain a statement that describes where the data supporting the conclusions can be obtained. Data cannot be listed as ''Available from authors'' or stored solely in supporting information. Citations to archived data should be included in your reference list. Wiley will publish it as a separate section on the paper’s page. Examples and complete information are here:
https://www.agu.org/Publish with AGU/Publish/Author Resources/Data for Authors


\acknowledgments
Enter acknowledgments here. This section is to acknowledge funding, thank colleagues, enter any secondary affiliations, and so on.


%%%%%%%%%%%%%%%%%%%%%%%%%%%%%%%%%%%%%%%%%%%%%%%
% REFERENCES and BIBLIOGRAPHY
%
\bibliography{refs} %don't specify the file extension
% don't specify bibliographystyle
%
%%%%%%%%%%%%%%%%%%%%%%%%%%%%%%%%%%%%%%%%%%%%%%%

%\bibliography{ enter your bibtex bibliography filename here }

% latex table generated in R 4.4.0 by xtable 1.8-4 package
% Thu Aug  1 12:17:48 2024
\begin{table}[ht]
\centering
\begin{tabular}{rllrrrrrr}
  \hline
 & basin & forcing\_type & month & cv\_1 & cv\_2 & cv\_3 & cv\_4 & por \\ 
  \hline
1 & FSSO3 & FA &    1 & 0.871 & 0.890 & 0.860 & 0.873 & 0.878 \\ 
  2 & FSSO3 & FA &    2 & 0.953 & 0.926 & 0.951 & 0.939 & 0.941 \\ 
  3 & FSSO3 & FA &    3 & 0.887 & 0.840 & 0.858 & 0.896 & 0.870 \\ 
  4 & FSSO3 & FA &    4 & 0.830 & 0.941 & 0.851 & 0.808 & 0.854 \\ 
  5 & FSSO3 & FA &    5 & 0.895 & 0.910 & 0.915 & 0.890 & 0.916 \\ 
  6 & FSSO3 & FA &    6 & 0.855 & 0.862 & 0.860 & 0.849 & 0.870 \\ 
  7 & FSSO3 & FA &    7 & 0.790 & 0.775 & 0.807 & 0.617 & 0.798 \\ 
  8 & FSSO3 & FA &    8 & 0.656 & 0.611 & 0.712 & 0.703 & 0.707 \\ 
  9 & FSSO3 & FA &    9 & 0.597 & 0.730 & 0.629 & 0.514 & 0.632 \\ 
  10 & FSSO3 & FA &   10 & 0.759 & 0.851 & 0.765 & 0.778 & 0.801 \\ 
  11 & FSSO3 & FA &   11 & 0.845 & 0.869 & 0.889 & 0.858 & 0.870 \\ 
  12 & FSSO3 & FA &   12 & 0.938 & 0.927 & 0.925 & 0.920 & 0.929 \\ 
  13 & FSSO3 & noFA &    1 & 0.756 & 0.746 & 0.774 & 0.720 & 0.752 \\ 
  14 & FSSO3 & noFA &    2 & 0.649 & 0.602 & 0.641 & 0.567 & 0.613 \\ 
  15 & FSSO3 & noFA &    3 & 0.731 & 0.757 & 0.740 & 0.646 & 0.726 \\ 
  16 & FSSO3 & noFA &    4 & 0.789 & 0.685 & 0.814 & 0.757 & 0.783 \\ 
  17 & FSSO3 & noFA &    5 & 0.821 & 0.753 & 0.789 & 0.799 & 0.796 \\ 
  18 & FSSO3 & noFA &    6 & 0.660 & 0.673 & 0.700 & 0.815 & 0.715 \\ 
  19 & FSSO3 & noFA &    7 & 0.613 & 0.665 & 0.797 & 0.696 & 0.698 \\ 
  20 & FSSO3 & noFA &    8 & 0.465 & 0.514 & 0.603 & 0.540 & 0.533 \\ 
  21 & FSSO3 & noFA &    9 & 0.678 & 0.753 & 0.667 & 0.633 & 0.658 \\ 
  22 & FSSO3 & noFA &   10 & 0.787 & 0.769 & 0.748 & 0.749 & 0.758 \\ 
  23 & FSSO3 & noFA &   11 & 0.806 & 0.836 & 0.758 & 0.848 & 0.812 \\ 
  24 & FSSO3 & noFA &   12 & 0.632 & 0.673 & 0.662 & 0.650 & 0.648 \\ 
   \hline
\end{tabular}
\end{table}
% latex table generated in R 4.4.0 by xtable 1.8-4 package
% Thu Aug  1 12:17:48 2024
\begin{table}[ht]
\centering
\begin{tabular}{rllrrrrrr}
  \hline
 & basin & forcing\_type & month & cv\_1 & cv\_2 & cv\_3 & cv\_4 & por \\ 
  \hline
1 & SAKW1 & FA &    1 & 0.839 & 0.819 & 0.871 & 0.867 & 0.800 \\ 
  2 & SAKW1 & FA &    2 & 0.821 & 0.808 & 0.778 & 0.738 & 0.770 \\ 
  3 & SAKW1 & FA &    3 & 0.680 & 0.717 & 0.698 & 0.685 & 0.682 \\ 
  4 & SAKW1 & FA &    4 & 0.771 & 0.774 & 0.715 & 0.730 & 0.739 \\ 
  5 & SAKW1 & FA &    5 & 0.664 & 0.680 & 0.767 & 0.644 & 0.701 \\ 
  6 & SAKW1 & FA &    6 & 0.788 & 0.840 & 0.814 & 0.783 & 0.794 \\ 
  7 & SAKW1 & FA &    7 & 0.876 & 0.879 & 0.866 & 0.842 & 0.885 \\ 
  8 & SAKW1 & FA &    8 & 0.650 & 0.793 & 0.626 & 0.680 & 0.653 \\ 
  9 & SAKW1 & FA &    9 & 0.743 & 0.844 & 0.692 & 0.838 & 0.775 \\ 
  10 & SAKW1 & FA &   10 & 0.898 & 0.933 & 0.880 & 0.866 & 0.896 \\ 
  11 & SAKW1 & FA &   11 & 0.809 & 0.804 & 0.802 & 0.805 & 0.846 \\ 
  12 & SAKW1 & FA &   12 & 0.851 & 0.796 & 0.867 & 0.769 & 0.871 \\ 
  13 & SAKW1 & noFA &    1 & 0.819 & 0.783 & 0.860 & 0.761 & 0.836 \\ 
  14 & SAKW1 & noFA &    2 & 0.805 & 0.808 & 0.755 & 0.709 & 0.793 \\ 
  15 & SAKW1 & noFA &    3 & 0.753 & 0.701 & 0.727 & 0.695 & 0.739 \\ 
  16 & SAKW1 & noFA &    4 & 0.659 & 0.635 & 0.630 & 0.636 & 0.647 \\ 
  17 & SAKW1 & noFA &    5 & 0.460 & 0.622 & 0.525 & 0.445 & 0.467 \\ 
  18 & SAKW1 & noFA &    6 & 0.750 & 0.823 & 0.771 & 0.718 & 0.745 \\ 
  19 & SAKW1 & noFA &    7 & 0.780 & 0.870 & 0.730 & 0.762 & 0.775 \\ 
  20 & SAKW1 & noFA &    8 & 0.744 & 0.875 & 0.672 & 0.671 & 0.730 \\ 
  21 & SAKW1 & noFA &    9 & 0.211 & 0.318 & 0.286 & 0.557 & 0.330 \\ 
  22 & SAKW1 & noFA &   10 & 0.781 & 0.766 & 0.823 & 0.786 & 0.767 \\ 
  23 & SAKW1 & noFA &   11 & 0.786 & 0.792 & 0.784 & 0.768 & 0.783 \\ 
  24 & SAKW1 & noFA &   12 & 0.817 & 0.744 & 0.797 & 0.701 & 0.795 \\ 
   \hline
\end{tabular}
\end{table}
% latex table generated in R 4.4.0 by xtable 1.8-4 package
% Thu Aug  1 12:17:48 2024
\begin{table}[ht]
\centering
\begin{tabular}{rllrrrrrr}
  \hline
 & basin & forcing\_type & month & cv\_1 & cv\_2 & cv\_3 & cv\_4 & por \\ 
  \hline
1 & WGCM8 & FA &    1 & 0.413 & 0.414 & 0.630 & 0.465 & 0.495 \\ 
  2 & WGCM8 & FA &    2 & 0.879 & 0.887 & 0.882 & 0.908 & 0.817 \\ 
  3 & WGCM8 & FA &    3 & 0.766 & 0.811 & 0.846 & 0.758 & 0.849 \\ 
  4 & WGCM8 & FA &    4 & 0.828 & 0.852 & 0.821 & 0.835 & 0.869 \\ 
  5 & WGCM8 & FA &    5 & 0.878 & 0.871 & 0.891 & 0.874 & 0.854 \\ 
  6 & WGCM8 & FA &    6 & 0.875 & 0.900 & 0.880 & 0.888 & 0.885 \\ 
  7 & WGCM8 & FA &    7 & 0.873 & 0.834 & 0.811 & 0.881 & 0.844 \\ 
  8 & WGCM8 & FA &    8 & 0.797 & 0.855 & 0.801 & 0.744 & 0.845 \\ 
  9 & WGCM8 & FA &    9 & 0.669 & 0.752 & 0.673 & 0.755 & 0.704 \\ 
  10 & WGCM8 & FA &   10 & 0.585 & 0.659 & 0.689 & 0.695 & 0.736 \\ 
  11 & WGCM8 & FA &   11 & 0.580 & 0.737 & 0.667 & 0.739 & 0.769 \\ 
  12 & WGCM8 & FA &   12 & 0.793 & 0.689 & 0.791 & 0.704 & 0.784 \\ 
  13 & WGCM8 & noFA &    1 & 0.478 & 0.379 & 0.640 & 0.515 & 0.470 \\ 
  14 & WGCM8 & noFA &    2 & 0.902 & 0.878 & 0.902 & 0.889 & 0.898 \\ 
  15 & WGCM8 & noFA &    3 & 0.802 & 0.799 & 0.832 & 0.734 & 0.779 \\ 
  16 & WGCM8 & noFA &    4 & 0.858 & 0.851 & 0.849 & 0.784 & 0.844 \\ 
  17 & WGCM8 & noFA &    5 & 0.845 & 0.848 & 0.818 & 0.840 & 0.844 \\ 
  18 & WGCM8 & noFA &    6 & 0.860 & 0.884 & 0.857 & 0.860 & 0.862 \\ 
  19 & WGCM8 & noFA &    7 & 0.906 & 0.853 & 0.836 & 0.888 & 0.879 \\ 
  20 & WGCM8 & noFA &    8 & 0.764 & 0.832 & 0.824 & 0.747 & 0.835 \\ 
  21 & WGCM8 & noFA &    9 & 0.580 & 0.581 & 0.461 & 0.581 & 0.488 \\ 
  22 & WGCM8 & noFA &   10 & 0.733 & 0.742 & 0.770 & 0.762 & 0.756 \\ 
  23 & WGCM8 & noFA &   11 & 0.671 & 0.702 & 0.705 & 0.644 & 0.694 \\ 
  24 & WGCM8 & noFA &   12 & 0.839 & 0.656 & 0.765 & 0.709 & 0.746 \\ 
   \hline
\end{tabular}
\end{table}

%Reference citation instructions and examples:
%
% Please use ONLY \cite and \citeA for reference citations.
% \cite for parenthetical references
% ...as shown in recent studies (Simpson et al., 2019)
% \citeA for in-text citations
% ...Simpson et al. (2019) have shown...
%
%
%...as shown by \citeA{jskilby}.
%...as shown by \citeA{lewin76}, \citeA{carson86}, \citeA{bartoldy02}, and \citeA{rinaldi03}.
%...has been shown \cite{jskilbye}.
%...has been shown \cite{lewin76,carson86,bartoldy02,rinaldi03}.
%... \cite <i.e.>[]{lewin76,carson86,bartoldy02,rinaldi03}.
%...has been shown by \cite <e.g.,>[and others]{lewin76}.
%
% apacite uses < > for prenotes and [ ] for postnotes
% DO NOT use other cite commands (e.g., \citet, \citep, \citeyear, \nocite, \citealp, etc.).
%



\end{document}



More Information and Advice:

%%%%%%%%%%%%%%%%%%%%%%%%%%%%%%%%%%%%%%%%%%%%%%%
%
%  SECTION HEADS
%
%%%%%%%%%%%%%%%%%%%%%%%%%%%%%%%%%%%%%%%%%%%%%%%

% Capitalize the first letter of each word (except for
% prepositions, conjunctions, and articles that are
% three or fewer letters).

% AGU follows standard outline style; therefore, there cannot be a section 1 without
% a section 2, or a section 2.3.1 without a section 2.3.2.
% Please make sure your section numbers are balanced.
% ---------------
% Level 1 head
%
% Use the \section{} command to identify level 1 heads;
% type the appropriate head wording between the curly
% brackets, as shown below.
%
%An example:
%\section{Level 1 Head: Introduction}
%
% ---------------
% Level 2 head
%
% Use the \subsection{} command to identify level 2 heads.
%An example:
%\subsection{Level 2 Head}
%
% ---------------
% Level 3 head
%
% Use the \subsubsection{} command to identify level 3 heads
%An example:
%\subsubsection{Level 3 Head}
%
%---------------
% Level 4 head
%
% Use the \subsubsubsection{} command to identify level 3 heads
% An example:
%\subsubsubsection{Level 4 Head} An example.
%
%%%%%%%%%%%%%%%%%%%%%%%%%%%%%%%%%%%%%%%%%%%%%%%
%
%  IN-TEXT LISTS
%
%%%%%%%%%%%%%%%%%%%%%%%%%%%%%%%%%%%%%%%%%%%%%%%
%
% Do not use bulleted lists; enumerated lists are okay.
% \begin{enumerate}
% \item
% \item
% \item
% \end{enumerate}
%
%%%%%%%%%%%%%%%%%%%%%%%%%%%%%%%%%%%%%%%%%%%%%%%
%
%  EQUATIONS
%
%%%%%%%%%%%%%%%%%%%%%%%%%%%%%%%%%%%%%%%%%%%%%%%

% Single-line equations are centered.
% Equation arrays will appear left-aligned.

Math coded inside display math mode \[ ...\]
 will not be numbered, e.g.,:
 \[ x^2=y^2 + z^2\]

 Math coded inside \begin{equation} and \end{equation} will
 be automatically numbered, e.g.,:
 \begin{equation}
 x^2=y^2 + z^2
 \end{equation}


% To create multiline equations, use the
% \begin{eqnarray} and \end{eqnarray} environment
% as demonstrated below.
\begin{eqnarray}
  x_{1} & = & (x - x_{0}) \cos \Theta \nonumber \\
        && + (y - y_{0}) \sin \Theta  \nonumber \\
  y_{1} & = & -(x - x_{0}) \sin \Theta \nonumber \\
        && + (y - y_{0}) \cos \Theta.
\end{eqnarray}

%If you don't want an equation number, use the star form:
%\begin{eqnarray*}...\end{eqnarray*}

% Break each line at a sign of operation
% (+, -, etc.) if possible, with the sign of operation
% on the new line.

% Indent second and subsequent lines to align with
% the first character following the equal sign on the
% first line.

% Use an \hspace{} command to insert horizontal space
% into your equation if necessary. Place an appropriate
% unit of measure between the curly braces, e.g.
% \hspace{1in}; you may have to experiment to achieve
% the correct amount of space.


%%%%%%%%%%%%%%%%%%%%%%%%%%%%%%%%%%%%%%%%%%%%%%%
%
%  EQUATION NUMBERING: COUNTER
%
%%%%%%%%%%%%%%%%%%%%%%%%%%%%%%%%%%%%%%%%%%%%%%%

% You may change equation numbering by resetting
% the equation counter or by explicitly numbering
% an equation.

% To explicitly number an equation, type \eqnum{}
% (with the desired number between the brackets)
% after the \begin{equation} or \begin{eqnarray}
% command. The \eqnum{} command will affect only
% the equation it appears with; LaTeX will number
% any equations appearing later in the manuscript
% according to the equation counter.
%

% If you have a multiline equation that needs only
% one equation number, use a \nonumber command in
% front of the double backslashes (\\) as shown in
% the multiline equation above.

% If you are using line numbers, remember to surround
% equations with \begin{linenomath*}...\end{linenomath*}

%  To add line numbers to lines in equations:
%  \begin{linenomath*}
%  \begin{equation}
%  \end{equation}
%  \end{linenomath*}



